\section{Problemstellung und Motivation}

Coding Agents sind in der professionellen Softwareentwicklung angekommen.
Coding Agents wie SWE-agent übernehmen nicht länger nur Autovervollständigung,
sondern planen Aufgaben, rufen Werkzeuge auf, lesen Dateien
und koordinieren Änderungen über mehrere Module hinweg~\parencite{yang2024sweagent}.
Kommerzielle Werkzeuge wie Claude Code, Cursor und GitHub Copilot Workspace
setzen dieses Paradigma mittlerweile im Produktiveinsatz um.
An der Literatur lässt sich ein messbarer Produktivitätsgewinn belegen:
Entwicklerinnen und Entwickler schließen abgegrenzte Aufgaben mit AI-Unterstützung
bis zu 55,8\,\% schneller ab als ohne \parencite{peng2023copilot}.
Der SWE-bench-Benchmark zeigt, dass selbst State-of-the-Art-Modelle
an realen Repository-Aufgaben nahezu vollständig scheitern:
Das beste getestete Modell (Claude~2) löst lediglich 1,96\,\% der Issues~\parencite{jimenez2024swebench},
weil die Aufgaben Multi-File-Koordination, lange Kontexte und komplexes Reasoning erfordern.

Ein wesentliches Hindernis liegt dabei nicht allein im Modell,
sondern in der Kontextbereitstellung:
Was ein Coding Agent über eine Codebase weiß
(welche Dateien er gelesen hat, welche Abhängigkeiten er kennt,
welche Konventionen ihm vertraut sind)
bestimmt entscheidend, ob er eine Aufgabe korrekt löst
oder in produktionsfremden Mustern halluziniert~\parencite{li2026contextbench}.

\subsection*{Das Reaktivitätsproblem moderner Coding Agents}

Aktuelle Coding Agents explorieren Codebases reaktiv.
Das bedeutet, dass der Agent eigenständig Dateien öffnet und nach relevanten Funktionen sucht,
wenn er eine konkrete Aufgabe erhält.
Diese Strategie ist teuer und fehleranfällig.
\Textcite{li2026contextbench} dokumentieren, dass Coding Agents systematisch
mehr Kontext abrufen als sie tatsächlich nutzen,
und dabei vorher besuchte Dateien wiederholt aufsuchen
(ein Muster redundanter Exploration).
\Textcite{suri2026codescout} zeigen, dass Agents bei unklaren Aufgabenbeschreibungen
in eine ziellose Überexploration geraten.
Ein schlecht formulierter Bug-Report führt in ihren Experimenten
zu 21 Suchschritten ohne Ergebnis.
Mit einer strukturierten Aufgabenbeschreibung löst derselbe Agent
das gleiche Problem in sechs Schritten.

Es gibt bereits proaktive Ansätze, die Kontext vor der Agentausführung aufbereiten.
\textcite{xia2024agentless} lokalisiert Fehlerstellen in einem vorgelagerten Dreiphasen-Prozess.
\textcite{suri2026codescout} reichert die Aufgabenbeschreibung
durch eine Vorab-Exploration des Repositories an.
Beide Systeme analysieren die Codebase aufgabenbezogen
und zum Zeitpunkt der jeweiligen Anfrage.
Eine fortlaufende Synchronisation mit einer aktiv weiterentwickelten Codebase
ist nicht vorgesehen.

\subsection*{Das Transparenzproblem}

Neben dem Reaktivitätsproblem stellt sich die Frage der Inspizierbarkeit.
Wenn ein Coding Agent eine falsche Lösung produziert,
kann der Entwickler nicht nachvollziehen,
welcher Kontext zu welcher Entscheidung geführt hat.
Kontext ist eine Black Box.
Fehlt einer Funktion die Kenntnis ihrer Aufrufer?
Hat der Agent das zugehörige Schema nicht gelesen?
Ohne Antwort auf diese Fragen bleibt die Fehlersuche schwer nachvollziehbar.
Erste Ansätze zur Agenten-Transparenz, etwa AgentStepper~\parencite{agentstepper2026},
helfen Entwicklern, Ausführungstrajektorien nachzuvollziehen.
Die Inspizierbarkeit der Kontextauswahl selbst
findet dabei bislang kaum Beachtung.

\subsection*{Zielsetzung}

Ziel dieser Arbeit ist die Entwicklung und Evaluation eines Living Knowledge Graph
als Kontextinfrastruktur für Coding Agents.
Das System repräsentiert eine Codebase als strukturierten,
automatisch synchronisierten Graphen aus Dateien, Funktionen, Abhängigkeiten
und Testergebnissen.
Erhält das System eine Aufgabe, sucht es anhand des Living Knowledge Graph selbstständig die passenden Code-Abschnitte heraus
und gibt sie dem Agenten weiter.
Es erklärt dabei, warum jede Stelle ausgewählt wurde.
Das System lässt sich direkt in bestehende Werkzeuge wie Claude Code und Cursor einbinden.

Die Evaluation erfolgt an laufenden Hochschulprojekten,
um den Mehrwert im realen Entwicklungsalltag zu messen.
Diese Arbeit soll zeigen, dass Coding Agents durch den Living Knowledge Graph
weniger Zeit damit verbringen, sich in einer Codebase zu orientieren,
und weniger unnötige Dateien in ihren Kontext aufnehmen.
Außerdem sollen Entwickler besser nachvollziehen können,
warum ein Agent eine bestimmte Lösung vorschlägt.
